% Risk--coverage, plotted from research/refusal_signal/results/baseline.jsonl.
% Every coordinate is generated by make_data.py; none is transcribed.
\begin{tikzpicture}
  \begin{axis}[reportaxis,
      xlabel={Coverage --- fraction of cases answered},
      ylabel={Risk --- error rate on answered set (\%)},
      xmin=0, xmax=1.02, ymin=-0.6, ymax=20,
      xtick={0,0.2,0.4,0.6,0.8,1.0},
      legend style={at={(0.97,0.06)}, anchor=south east},
      legend cell align=left]

    % answer-everything reference
    \addplot[draw=slate, dashed, line width=.8pt, forget plot, mark=none]
      coordinates {(0,\rcAnswerAllRisk) (1.02,\rcAnswerAllRisk)};
    \node[font=\plexsans\tiny, text=slate, anchor=south west]
      at (axis cs:0.42,\rcAnswerAllRisk) {answer everything, \rcAnswerAllRisk\,\%};

    % the confidence selector
    \addplot[draw=forge, line width=1.2pt, mark=none]
      table[col sep=comma, x=coverage, y=risk] {data/rc_confidence.csv};
    \addlegendentry{self-reported confidence}

    % the refusal selector: one operating point, not a curve
    \addplot[only marks, mark=*, mark size=2.6pt, draw=rosedeep, fill=rosedeep]
      table[col sep=comma, x=coverage, y=risk] {data/rc_refusal_point.csv};
    \addlegendentry{abstain unless IR-covered}

    \node[font=\plexsans\tiny, text=rosedeep, anchor=west]
      at (axis cs:0.17,16.9) {worse than answering everything};
  \end{axis}
\end{tikzpicture}
